% !TEX TS-program = xelatex
% !TEX encoding = UTF-8 Unicode
% Migrated from 正式简历-草稿版.md; edit the content below directly.
\documentclass{resume}
\usepackage{zh_CN-Adobefonts_external}
\usepackage{linespacing_fix}
\geometry{left=1.8cm,right=1.8cm,top=1.5cm,bottom=1.5cm}
\definecolor{resumegray}{gray}{0.35}
\titleformat{\section}{\fontsize{13pt}{16pt}\selectfont\bfseries\raggedright}{}{0em}{}[\titlerule]
\titlespacing*{\section}{0pt}{8pt}{5pt}
\titleformat{\subsection}{\normalsize\bfseries\raggedright}{}{0em}{}
\titlespacing*{\subsection}{0pt}{6pt}{3pt}
\setlist[itemize]{leftmargin=1.2em,itemsep=2pt,parsep=0pt,topsep=3pt,partopsep=0pt}
\setlength{\emergencystretch}{3em}
\newcommand{\entry}[3]{%
  \subsection*{#1}%
  {\small\color{resumegray}#2\hfill #3\par}\nopagebreak[4]
}
\newcommand{\module}[1]{\par\vspace{5pt}\textbf{#1}\par\nopagebreak[4]}
\renewcommand{\name}[1]{%
  \centerline{{\fontsize{22pt}{26pt}\selectfont\scshape #1}}%
  \vspace{1.2ex}%
}
\begin{document}
\fontsize{10.5pt}{13.5pt}\selectfont
\pagenumbering{gobble}
\name{翟开放}
{\centering 15529607820\quad\textperiodcentered\quad zkfopen@foxmail.com\quad\textperiodcentered\quad 1995 年生\quad\textperiodcentered\quad 北京\par}
\vspace{4pt}



\section{职业概述}

6 年以上后端、云原生与 AI Infra 研发经验。曾在百度从事 Kubernetes 私有化交付平台研发，现于智谱负责 MaaS 平台研发与模型部署，重点深入昇腾平台的推理性能优化。具备 PD 分离、长上下文与 KV Cache 验证经验，能够结合性能测试开展部署选型与容量规划。\par

\section{技术能力}

\begin{itemize}

  \item \textbf{模型部署：}vLLM-Ascend、SGLang、MindIE；昇腾 910B/910C/950PR；W8A8 等量化模型部署、PD 分离、DP/TP/EP。

  \item \textbf{推理性能：}EvalScope perf、TTFT/TPOT、吞吐与容量分析；长上下文、前缀缓存、CPU KV 回载、UCM。

  \item \textbf{云原生与 MaaS：}Linux（Ubuntu、CentOS、麒麟）、Docker、Kubernetes、ARM64/AMD64、异构算力纳管、资源池与模型服务；Helm、Argo CD、CI/CD。

  \item \textbf{后端与自动化：}Go、client-go、Gin、Gorm、Python、Shell；权限与审计、REST API、MCP 服务。

  \item \textbf{可观测性与排障：}Prometheus、Grafana、Loki；模型服务、存储、网络、HCCL/RDMA 通信排障。

\end{itemize}

\section{工作经历}

\entry{智谱｜大模型开发工程师}{商业智能事业部}{2025.05—至今}

\begin{itemize}

  \item \textbf{模型部署与交付：}负责 GLM、DeepSeek、Qwen 等模型的部署适配与方案选型，以昇腾 910B/910C/950PR 为主，面向中关村实验室、山东移动、中石化等客户完成模型部署上线、推理调优及服务故障排查。

  \item \textbf{MaaS 平台研发：}主导资源池、模型用量统计与算力监控等核心模块建设，完成 ARM64 适配及通算/智算双集群部署，打通模型实例下发链路；完成模型调用向 NewAPI 的迁移，保障迁移期间业务连续性。

  \item \textbf{推理性能工程：}建设 EvalScope perf 自动压测，开展长上下文、PD 拓扑与 KV Cache 方案对比，明确不同负载下的容量边界，输出部署选型、并发配置建议及标准化交付包。

\end{itemize}

\entry{百度｜高级研发工程师}{百度智能云事业群}{2020.07—2025.05}

\begin{itemize}

  \item \textbf{云原生平台研发：}负责天牛 Kubernetes 集群生命周期与资源管理核心模块设计开发，建设权限校验、操作审计及多版本兼容能力，支撑 3000+ 集群累计托管与 1500+ 项目交付。

  \item \textbf{异构适配与制品交付：}完成 ARM64/AMD64、Ubuntu/麒麟环境适配，建设镜像及模型制品导入导出、差异资源交付与主辅包独立部署能力，支持异构环境下的私有化交付。

  \item \textbf{自动化与交付保障：}基于 Go、Kubernetes、Helm 建设巡检及集群预检查工具，完善代码评审、单元测试与自动化发布链路，将环境检查接入集群创建和扩容流程。

\end{itemize}

\section{重点项目}

\entry{中关村实验室 - 大模型 PD 分离部署与长上下文性能验证}{项目技术负责人}{2026.06—至今}

面向昇腾 910C 长上下文推理场景，开展 GLM-5.2、GLM-5.3、DeepSeek-V4-Pro 模型部署工程化、PD 拓扑选型与缓存方案评估。\par

\begin{itemize}

  \item \textbf{部署工程化：}在昇腾 910C 环境完成 GLM-5.2/5.3 量化模型的 PD 分离部署与 MaaS 上架，将 3P1D Docker 方案封装为 Helm Chart，通过 Argo CD 同步配置并验证扩缩容；完成 DeepSeek-V4-Pro 部署与问答验证。

  \item \textbf{拓扑选型：}基于 EvalScope perf 开展长上下文与多并发测试；在约 512K 输入、TTFT＜30 秒及单路 Decode≥15 token/s 的目标下，对比验证 2P1D、1P1D 的并发承载能力，分别约为 16、8 并发，形成对应场景的拓扑与容量建议。

  \item \textbf{缓存方案评估：}对比 UCM 与合作方 KV Cache 优化版本：UCM 在 972.8K/8 并发下由全量失败变为 64/64 成功，但多数档位吞吐下降，判定收益有限；合作方版本在 786K/973K 下 TTFT 最大 223.07→23.48 秒、整体吞吐最大 12,601→60,941 tok/s，据此给出方案适用范围。

  \item \textbf{排障与达标验证：}定位 Proxy 启动与请求分发异常，修复后由 1P1D 调整为 3P1D；GLM-5.3 在 1024K/4 并发/90\% 前缀下 TTFT 13.72 秒、TPOT 49.06ms，满足 $\leq$~40s/$\leq$~50ms 目标。

\end{itemize}

\newpage

\entry{山东移动 - 国产化 MaaS 与模型部署优化}{项目技术负责人}{2025.06—至今}

建设 ARM64 与昇腾环境下的私有化 MaaS 平台，支撑百卡级算力集群的统一管理与模型服务部署，推进 GLM-5.3-Flash 部署适配与性能验证。\par

\begin{itemize}

  \item \textbf{国产化 MaaS 建设：}完成 MaaS API、Infra 及配套组件的 ARM64 适配，建设通算、智算双 Kubernetes 集群并打通模型实例下发与调用链路；为 GLM、DeepSeek、Qwen 等 13 个去重模型制定部署方案并指导现场实施，将 TTS、ASR 服务封装为 Helm Chart 纳管至 MaaS。

  \item \textbf{Flash 部署与压测：}基于 vLLM-Ascend 制定 GLM-5.3-Flash-W8A8 的 910C 单机 TP16、910B 单机 TP8 及双机 1P1D 部署方案（P 节点：DP1/TP1/PCP8；D 节点：DP4/TP2/PCP1），实现部署脚本化与 EvalScope perf 自动压测，完成不同硬件与拓扑的性能对比。

  \item \textbf{KV Cache 验证：}在 910B、16K 前缀受控测试中，验证 CPU KV Offloading 使 TTFT 从 4.14 秒降至 0.822 秒；在 910B 双机、512K 共享前缀输入、8 并发下，测得单路 Decode 超过 10 token/s。

\end{itemize}

\entry{中石化 - 人工智能中台 MaaS 与模型服务}{项目技术负责人}{2025.05—2026.04}

面向大型能源央企建设私有化 MaaS 平台，统一管理千卡级异构算力与模型服务，支持模型推理、微调及 API 调用。\par

\begin{itemize}

  \item \textbf{模型与算力交付：}累计交付 50+ 个去重模型，主导 GLM-4.7、DeepSeek-V3 部署与性能验证；将近 200 台裸金属服务器纳管至 Kubernetes，参与业务与昇腾算力集群拆分，制定节点下架、模型服务及对象存储迁移方案。

  \item \textbf{生产升级与保障：}负责 MaaS v2.0→v2.2 升级及定制代码、数据库、脚本迁移，通过演练、临时模型服务、网关切换与回滚准备保障升级，完成 5 个模型推理、2 个模型训练验证，验收无阻塞问题；定位 HCCL/RDMA 多机通信及 Kubernetes 故障，沉淀部署运维 SOP。

\end{itemize}

\entry{百度天牛私有化交付平台}{核心模块负责人}{2020.07—2025.05}

面向私有化交付场景，长期负责集群管理、制品交付及自动化巡检模块的设计、研发与迭代。\par

\module{模块一：Kubernetes 集群生命周期与异构平台建设}

\begin{itemize}

  \item 设计集群生命周期框架，覆盖创建、扩缩容等 11 项功能，支持 Kubernetes 1.20/1.24/1.28；完成 ARM64/AMD64、Ubuntu/CentOS/麒麟适配，支撑 60+ 多架构集群稳定运行，完善权限校验、审计与异常处理，平台交付时效提升 40\%、故障率下降 65\%。

\end{itemize}

\module{模块二：大规模制品交付与存储优化}

\begin{itemize}

  \item 设计镜像架构识别与导出回退机制，重构镜像及模型文件加载逻辑，将 1TB 产品包导入由 2 小时缩短至 1 分钟，累计支撑 37,000+ 次导出；支持差异资源导出与主辅包独立部署，通过 CFS 存储与资源治理，将年度存储成本由约 40 万元降至约 1 万元。

\end{itemize}

\module{模块三：自动化巡检与集群预检查}

\begin{itemize}

  \item 建设覆盖 300+ 业务场景的自动化巡检与 CLI 工具，将 15 项预检查接入集群创建、扩容流程，以 goroutine 并发执行；累计完成 10,000+ 次前置检查，集群启动成功率由 80\% 提升至 98\%，年均节省运维人力 150+ 人天。

\end{itemize}

\section{荣誉认证}

\begin{itemize}

  \item 发明专利 2 项（2022 年），其中《资源信息处理方法、装置、电子设备以及存储介质》为第一作者。

  \item 获百度 ACG 体系交付服务奖（2023 Q2）、ACG 靠谱工程师（2021 Q4）及百度小赞奖。

  \item Kubernetes CKA 认证｜2025.01

\end{itemize}

\section{教育经历}

{\renewcommand{\arraystretch}{1.35}
\begin{tabular*}{\textwidth}{@{}l@{\extracolsep{\fill}}llr@{}}
西安电子科技大学 & 计算机科学与技术 & 硕士学位 & 2017.09—2020.06 \\
河南工业大学 & 软件工程 & 学士学位 & 2013.09—2017.06 \\
\end{tabular*}\par}

\end{document}
